\chapter{Pemrograman Linear dan Integer}

Sebuah bengkel kayu menerima 24 lembar kayu lapis dan 6 jam mesin perakitan untuk hari
itu. Satu rak buku memakai 6 lembar kayu lapis dan 1 jam perakitan. Satu bangku memakai
4 lembar dan 2 jam. Rak memberi laba 5 satuan, bangku 4 satuan. Berapa banyak rak dan
bangku yang sebaiknya dibuat?

Naluri pertama memilih rak, sebab labanya per unit lebih tinggi. Naluri itu salah: satu rak
menghabiskan kayu lapis sebanyak satu setengah bangku. Jawaban terbaik model bersambungnya
3 rak dan 1,5 bangku, berlaba 21, dan setengah bangku tidak bisa dijual. Pembulatan ke bawah
memberi 19; 4 rak tanpa bangku memberi 20. Jadi jawaban bulat terbaik bukan pembulatan
jawaban pecahan terbaik, dan selisihnya kesalahan cara berpikir, bukan kesalahan hitung.

Alat dari bab sebelumnya belum cukup. Bab 8 mengajari kita menyelesaikan
$A\mathbf{x}=\mathbf{b}$, tetapi di sini persamaannya lebih sedikit daripada variabelnya,
jadi ada tak berhingga jawaban dan kita harus memilih yang terbaik. Bab 23 memberi kondisi
KKT, tetapi pada soal ini gradien sasaran konstan dan tidak pernah lenyap: jawabannya
menempel di batas daerah layak, bukan di titik stasioner. Jaminan bab 21 bahwa optimum lokal
juga global tidak menunjukkan di mana optimumnya, dan hilang begitu variabel dituntut
bulat.

Karena itu bab ini menempuh tiga langkah: menulis soal sebagai program linear dan
mengenali bahwa jawabannya ada di sudut daerah layak; memberi harga pada setiap kendala
lewat dualitas, sehingga keoptimalan dapat diperiksa tanpa menyisir seluruh daerah; lalu
menghadapi tuntutan kebulatan, saat model bersambung memberi batas yang sah tetapi belum
tentu tercapai.

\begin{untukanak}
Bayangkan kamu punya 24 lembar kertas warna dan 6 jam untuk membuat dua jenis hiasan
kelas. Hiasan pertama memakai 6 lembar dan 1 jam, hiasan kedua memakai 4 lembar dan 2 jam.
Hiasan pertama lebih banyak dipesan, jadi harganya lebih tinggi. Berapa banyak masing-masing
yang sebaiknya kamu buat?

Kalau kamu coba-coba, ada banyak campuran yang muat. Yang mengejutkan, campuran paling
menguntungkan ternyata memakai setengah hiasan. Setengah hiasan tidak bisa dijual, dan di
sinilah akal sehat mulai menipu: kalau kamu membulatkan ke bawah, hasilnya bukan yang
terbaik. Ada campuran lain yang seluruhnya bulat dan justru lebih menguntungkan, padahal
tidak kelihatan seperti pembulatan apa pun.

Pertanyaan yang membuat penasaran begini. Kalau ada ribuan campuran yang mungkin, bagaimana
kita tahu satu campuran itu yang terbaik tanpa mencoba semuanya? Dan kalau temanmu
menawarkan satu lembar kertas tambahan, berapa harga wajar untuk lembar itu? Di bab ini
kita akan tahu mengapa jawaban terbaik selalu ada di pojok daerah yang mungkin, mengapa
setiap bahan punya harga yang bisa dihitung, dan mengapa jawaban pecahan harus dibulatkan
dengan hati-hati, bukan dengan naluri.
\end{untukanak}

\section{Dari rencana produksi ke bentuk baku}
\tingkat{dasar}

Satu rencana produksi adalah sepasang bilangan: banyak rak dan banyak bangku. Rencana
disebut layak bila pemakaian bahannya tidak melebihi persediaan. Untuk rak $x_{1}$ dan
bangku $x_{2}$, syaratnya
\[
  6x_{1}+4x_{2}\le24,\qquad x_{1}+2x_{2}\le6,\qquad x_{1}\ge0,\quad x_{2}\ge0 .
\]
Tiap syarat memotong bidang dengan satu garis dan menyisakan satu sisi; irisannya segiempat
dengan empat sudut. Labanya $5x_{1}+4x_{2}$, dan kita ingin angka itu sebesar mungkin.

Pola soal semacam ini disebut program linear. Namanya tepat: semua yang terlibat berbentuk
garis lurus, tidak ada lengkungan, tidak ada pangkat.

\kamus{program linear}{linear program}{Soal mencari nilai terbaik dari sesuatu yang dihitung
dengan penjumlahan berbobot, ketika semua aturannya berbentuk garis lurus.}

\begin{definisi}[Program linear]
Program linear adalah masalah memaksimumkan atau meminimumkan fungsi linear
$\mathbf{c}^{\top}\mathbf{x}$ terhadap variabel $\mathbf{x}\in\mathbb{R}^{n}$ dengan
kendala berbentuk ketaksamaan atau kesamaan linear. Himpunan titik yang memenuhi semua
kendala disebut \emph{daerah layak}. Bila daerah layaknya kosong, program disebut
\emph{tidak layak}; bila nilai sasarannya dapat dibuat sebesar apa pun, program disebut
\emph{tak terbatas}.\index{program linear}
\end{definisi}

Perangkat lunak dan buku teks memakai satu pola tulisan yang seragam, dan pola itulah yang
dipakai sepanjang bab ini.

\begin{definisi}[Bentuk baku]
Program linear dalam \emph{bentuk baku} adalah
\[
  \text{maksimumkan } \mathbf{c}^{\top}\mathbf{x}
  \quad\text{dengan}\quad A\mathbf{x}=\mathbf{b},\qquad \mathbf{x}\ge\mathbf{0},
\]
dengan $A$ berukuran $m\times n$ berperingkat $m$, $\mathbf{b}\in\mathbb{R}^{m}$, dan
$\mathbf{c}\in\mathbb{R}^{n}$. Setiap program linear dapat dibawa ke bentuk ini tanpa
mengubah nilai optimumnya.\index{bentuk baku}
\end{definisi}

Empat terjemahan melakukan seluruh pekerjaan itu. Kendala
$\mathbf{a}^{\top}\mathbf{x}\le b$ menjadi kesamaan dengan menambah variabel yang menampung
sisa pemakaian; kendala $\mathbf{a}^{\top}\mathbf{x}\ge b$ dengan mengurangkannya. Variabel
yang boleh bernilai negatif diganti selisih dua variabel tak negatif, dan memaksimumkan
$\mathbf{c}^{\top}\mathbf{x}$ sama dengan meminimumkan $-\mathbf{c}^{\top}\mathbf{x}$.

\kamus{variabel kendur}{slack variable}{Bilangan yang dicatat sebagai sisa pemakaian, supaya
``tidak boleh lebih dari'' berubah menjadi ``sama dengan''.}

Untuk bengkel tadi, bentuk bakunya memakai $s_{1}$ sebagai sisa kayu lapis dan $s_{2}$
sebagai sisa jam perakitan:
\[
  \begin{array}{rcl}
    6x_{1}+4x_{2}+s_{1} &=& 24,\\
    x_{1}+2x_{2}+s_{2} &=& 6,
  \end{array}
  \qquad x_{1},x_{2},s_{1},s_{2}\ge0 .
\]
Di sini $m=2$ dan $n=4$, sehingga $A$ berukuran $2\times4$. Kedua variabel kendur itu bukan
hiasan: nanti keduanya muncul sebagai kolom di tabel simpleks dan sebagai petunjuk kendala
mana yang masih longgar.

Bentuk baku juga menuntut $\mathbf{b}\ge\mathbf{0}$, dan itu selalu bisa diatur
dengan mengalikan satu baris dengan $-1$. Yang tidak bisa diatur adalah peringkat: baris $A$
yang bergantung linear harus dibuang lebih dulu, sebab kalau tidak, matriks basis yang dipilih
bisa singular.

\section{Geometri: mengapa jawaban ada di sudut}
\tingkat{menengah}

Empat sudut segiempat tadi adalah $(0;0)$, $(4;0)$, $(3;1{,}5)$, dan $(0;3)$, dengan laba
$0$, $20$, $21$, dan $12$. Yang terbaik ada di sudut, dan itu bukan kebetulan: pada program
linear nilai terbaik selalu dapat dicari di sudut, sebab itulah metode simpleks hanya
berjalan dari sudut ke sudut.

\begin{figure}[!ht]
\centering
\begin{tikzpicture}
\begin{axis}[width=9.3cm,height=6.3cm,
  xlabel={$x_{1}$ (rak)},ylabel={$x_{2}$ (bangku)},
  xmin=-0.35,xmax=5.1,ymin=-0.35,ymax=3.9,
  axis lines=left,xtick={0,1,2,3,4,5},ytick={0,1,2,3},
  tick label style={font=\scriptsize},label style={font=\scriptsize}]
\addplot[forget plot,fill=layerteal!10,draw=layerteal,thick]
  coordinates {(0,0) (4,0) (3,1.5) (0,3)} \closedcycle;
\addplot[forget plot,domain=0:2.4,densely dotted,thick,layerindigo] {3-1.25*x};
\addplot[forget plot,domain=0:3.6,densely dotted,thick,layerindigo] {4.5-1.25*x};
\addplot[forget plot,domain=0.2:4.2,thick,accent] {5.25-1.25*x};
\addplot[forget plot,only marks,mark=*,mark size=1.7pt,draw=ink]
  coordinates {(0,0) (4,0) (0,3)};
\addplot[forget plot,only marks,mark=*,mark size=2.6pt,draw=accent]
  coordinates {(3,1.5)};
\node[font=\scriptsize,anchor=north west] at (axis cs:0,0) {$(0;0)$};
\node[font=\scriptsize,anchor=north east] at (axis cs:4,0) {$(4;0)$};
\node[font=\scriptsize,anchor=south west] at (axis cs:0,3) {$(0;3)$};
\node[font=\scriptsize,anchor=west] at (axis cs:3.05,1.45) {$(3;1{,}5)$};
\node[font=\scriptsize,anchor=west,layerindigo] at (axis cs:0.15,1.15) {$z=12$};
\node[font=\scriptsize,anchor=west,layerindigo] at (axis cs:0.15,2.15) {$z=18$};
\node[font=\scriptsize,anchor=west,accent] at (axis cs:0.15,3.45) {$z=21$};
\draw[-{Stealth},thick,accent] (axis cs:2.3,2.75) -- (axis cs:3.25,3.35)
  node[font=\scriptsize,above,accent,pos=0.9] {laba naik};
\end{axis}
\end{tikzpicture}
\caption{Daerah layak bengkel dan tiga garis laba sama besar. Poligon punya empat verteks,
$(0;0)$, $(4;0)$, $(3;1{,}5)$, $(0;3)$, bernilai $0$, $20$, $21$, $12$. Garis $z=12$ dan
$z=18$ masih memotong daerah layak; garis $z=21$ hanya menyentuh verteks $(3;1{,}5)$,
sehingga tidak ada titik layak yang bernilai lebih tinggi.}
\label{fig:ch22-daerah-optimum}
\end{figure}

Gambar~\ref{fig:ch22-daerah-optimum} mempertajam kebetulan itu. Tiga garis di sana adalah
garis laba sama besar, $5x_{1}+4x_{2}=12$, $18$, dan $21$. Garis laba yang lebih tinggi
lebih jauh dari titik nol, dan menggesernya ke arah itu memindahkannya keluar daerah layak;
garis $21$ hanya menyentuh verteks $(3;1{,}5)$.

Sudut punya nama yang lebih tepat, dan namanya tidak mengacu pada gambar.

\begin{definisi}[Verteks]
Misalkan $P\subseteq\mathbb{R}^{n}$ himpunan konveks. Titik $\mathbf{x}\in P$ disebut
\emph{verteks} atau \emph{titik ekstrem} bila tidak ada $\mathbf{y},\mathbf{z}\in P$ dengan
$\mathbf{y}\ne\mathbf{z}$ dan $t\in(0,1)$ sehingga
$\mathbf{x}=t\mathbf{y}+(1-t)\mathbf{z}$. Dengan kata lain, $\mathbf{x}$ bukan titik
tengah ruas antara dua titik lain di $P$.\index{verteks}
\end{definisi}

Pada segiempat tadi, titik $(3;1{,}5)$ adalah verteks: setiap ruas di dalam segiempat yang
melewatinya harus berakhir di titik itu sendiri. Sebaliknya $(2;1)$ bukan verteks, sebab ia
titik tengah antara $(1;1)$ dan $(3;1)$.

\kamus{verteks}{vertex}{Sudut daerah layak, yaitu titik yang tidak pernah menjadi titik
tengah antara dua titik layak lain.}

Langkah pertama pembuktian adalah mengubah penyangkalan menjadi arah yang bisa dipakai.

\begin{lemma}[Titik bukan verteks memberi garis]
Bila $\mathbf{x}\in P$ bukan verteks, maka ada $\mathbf{d}\ne\mathbf{0}$ sehingga
$\mathbf{x}+\mathbf{d}\in P$ dan $\mathbf{x}-\mathbf{d}\in P$.
\end{lemma}

\begin{proof}
Karena $\mathbf{x}$ bukan verteks, ada $\mathbf{y}\ne\mathbf{z}\in P$ dan $t\in(0,1)$
dengan $\mathbf{x}=t\mathbf{y}+(1-t)\mathbf{z}$. Tetapkan
$s=\tfrac12\min\{t,1-t\}>0$ dan $\mathbf{d}=s(\mathbf{y}-\mathbf{z})$. Maka
$\mathbf{x}+\mathbf{d}=(t+s)\mathbf{y}+(1-t-s)\mathbf{z}$ dan
$\mathbf{x}-\mathbf{d}=(t-s)\mathbf{y}+(1-t+s)\mathbf{z}$. Karena $0<t\pm s<1$, kedua
titik itu kombinasi konveks $\mathbf{y}$ dan $\mathbf{z}$, jadi anggota $P$ sebab $P$
konveks. Juga $\mathbf{d}\ne\mathbf{0}$, sebab $\mathbf{y}\ne\mathbf{z}$ dan $s>0$.
\end{proof}

Sekarang pernyataan yang menjadi alasan metode simpleks masuk akal. Kalimatnya dalam
bahasa sehari-hari: kalau daerah layaknya terkurung dan sasarannya tidak tumbuh ke arah
mana pun, maka puncak tertinggi ada di salah satu sudut, bukan di tengah bidang.

\begin{teorema}[Optimum tercapai di verteks]
Misalkan $P=\{\mathbf{x}:A\mathbf{x}\le\mathbf{b}\}$ tak kosong dan terbatas, dan misalkan
$\mathbf{c}^{\top}\mathbf{x}$ terbatas di atas pada $P$. Maka ada verteks
$\mathbf{x}^{*}$ dari $P$ dengan $\mathbf{c}^{\top}\mathbf{x}^{*}
=\max_{\mathbf{x}\in P}\mathbf{c}^{\top}\mathbf{x}$.
\index{optimum di verteks}
\end{teorema}

\begin{proof}
Himpunan optimal $F=\{\mathbf{x}\in P:\mathbf{c}^{\top}\mathbf{x}=p^{*}\}$ tak kosong,
tertutup, terbatas, dan konveks, jadi $F$ juga polytope. Tulis kendala $P$ sebagai
$\mathbf{a}_{i}^{\top}\mathbf{x}\le b_{i}$, $i=1,\dots,k$. Pilih $\mathbf{x}^{*}\in F$
yang memaksimumkan banyaknya kendala yang aktif, yaitu kendala dengan
$\mathbf{a}_{i}^{\top}\mathbf{x}^{*}=b_{i}$.

Andaikan $\mathbf{x}^{*}$ bukan verteks. Menurut lemma, ada $\mathbf{d}\ne\mathbf{0}$
dengan $\mathbf{x}^{*}\pm\mathbf{d}\in F$. Setiap kendala yang aktif di
$\mathbf{x}^{*}$ juga aktif di kedua titik itu, sehingga
$\mathbf{a}_{i}^{\top}\mathbf{d}=0$ bagi semua $i$ aktif. Tinjau garis
$\mathbf{x}^{*}+s\mathbf{d}$ untuk $s\in\mathbb{R}$. Irisannya dengan $P$ adalah selang
tertutup $[s_{-},s_{+}]$ yang memuat $[-1,1]$, dan karena $P$ terbatas kedua ujungnya
berhingga. Ujung $s_{+}$ ditentukan oleh kendala dengan $\mathbf{a}_{i}^{\top}\mathbf{d}>0$ yang
paling cepat tercapai; kendala itu aktif di $s_{+}$ tetapi tidak di $\mathbf{x}^{*}$, jadi
titik baru itu layak dan aktif di lebih banyak kendala. Nilainya tetap $p^{*}$ sebab
$\mathbf{c}^{\top}(\mathbf{x}^{*}+s_{+}\mathbf{d})=p^{*}+s_{+}\mathbf{c}^{\top}
\mathbf{d}=p^{*}$, dan $\mathbf{c}^{\top}\mathbf{d}=0$ karena kedua titik
$\mathbf{x}^{*}\pm\mathbf{d}$ bernilai $p^{*}$. Titik baru itu anggota $F$ dengan lebih banyak kendala
aktif, bertentangan dengan pilihan $\mathbf{x}^{*}$; maka $\mathbf{x}^{*}$ verteks.
\end{proof}

\begin{bacakeras}
Dua suara, dibaca bergantian dengan teman. Ketuk meja sekali setiap kali kata ``puncak''
diucapkan.

``Di mana jawaban terbaik?'' ``Di puncak.''
``Kenapa bukan di tengah?'' ``Karena di tengah ada tetangga di dua sisi.''
``Kalau ada tetangga, kenapa salah?'' ``Karena masih ada arah untuk bergeser.''
``Bergeser ke arah mana?'' ``Ke arah yang menanjak.''
``Kalau tidak ada arah menanjak?'' ``Maka kita sudah berdiri di puncak.''

Ulangi tiga baris tengah, makin pelan, dengan ketukan yang sama: di tengah ada tetangga, di
tengah masih bisa bergeser, puncak tidak punya arah menanjak.

Baris yang harus diulang pembaca, sekali, dengan suara jelas: jawaban terbaik ada di puncak,
bukan di tengah, dan puncak dikenali dari tidak adanya arah yang menanjak.
\end{bacakeras}

Bukti itu memakai keterbatasan $P$. Bila $P$ tak terbatas tetapi sasarannya
terbatas di atas, kesimpulannya masih benar; jalannya lewat dekomposisi $P$ menjadi selubung
konveks verteks ditambah kerucut arah resesi, dan bukti lengkapnya ada pada
\texttt{WSPrinted.pdf}, bagian struktur himpunan polihedral.

Banyaknya verteks berhingga, tetapi tidak selalu sedikit. Dalam bentuk baku,
verteks bersesuaian dengan pemilihan $m$ kolom $A$ yang bebas linear: kolom terpilih
menentukan $\mathbf{x}_{B}=B^{-1}\mathbf{b}$, kolom lain nol, dan titik itu layak bila
$\mathbf{x}_{B}\ge\mathbf{0}$. Calonnya sebanyak $\binom{n}{m}$, yang untuk $n=100$ dan
$m=50$ melampaui $10^{29}$; memeriksa semuanya bukan rencana masuk akal. Di sinilah simpleks
berperan: ia berjalan dari verteks ke verteks tetangga yang lebih baik.

\begin{tangan}
Bahan: lantai berpetak, dua utas tali atau selotip kertas, dan segenggam kancing.

Bentangkan dua tali lurus sebagai dua batas; satu untuk kayu lapis, satu untuk jam
perakitan. Berdirilah di petak pojok dan sebut petak itu titik nol: setiap petak ke kanan
menambah satu rak, setiap petak ke depan menambah satu bangku.

Berjalanlah menyusuri tali dan berhenti hanya di petak tempat dua tali bertemu. Letakkan
kancing di situ, dan di situ saja. Yang harus terlihat: semua kancing menempel di tikungan
pagar, tidak satu pun berdiri di tengah halaman. Sesudah itu hitung laba di setiap kancing
dengan jari, lalu pindahkan satu kancing ke tikungan yang paling menguntungkan.

Geser satu tali satu petak ke dalam, lalu ulangi: kancingnya berpindah, pemenangnya tetap
kancing di tikungan.
\end{tangan}

\section{Metode simpleks: satu tabel, dua pivot}
\tingkat{menengah}

Gagasannya sederhana. Berdiri di satu verteks, periksa apakah ada arah sisi yang menaikkan
laba; bila ada, berjalan ke verteks di ujungnya. Bila tidak, berhenti: tidak ada tetangga
yang lebih baik, yang pada program linear berarti tidak ada titik lain yang lebih baik.

Arah sisi itu ditentukan oleh \emph{biaya tereduksi}: seberapa banyak laba bertambah bila
satu variabel yang sekarang nol dinaikkan satu satuan, dengan mengorbankan variabel basis
secara optimal.

\kamus{solusi layak basis}{basic feasible solution}{Titik yang diperoleh dengan menyetel
sebagian variabel ke nol lalu menyelesaikan sisanya, dan kebetulan masih memenuhi semua
batas; titik ini selalu verteks.}

Untuk bengkel tadi, mulailah dari verteks $(0;0)$, yaitu $s_{1}=24$ dan $s_{2}=6$. Tabel
pertamanya memuat baris laba $z$ dan satu baris untuk tiap kendala:
\begin{center}
\begin{tabular}{c|cccc|c}
\hline
basis & $x_{1}$ & $x_{2}$ & $s_{1}$ & $s_{2}$ & ruas kanan \\
\hline
$z$   & $-5$ & $-4$ & $0$ & $0$ & $0$ \\
$s_{1}$ & $6$ & $4$ & $1$ & $0$ & $24$ \\
$s_{2}$ & $1$ & $2$ & $0$ & $1$ & $6$ \\
\hline
\end{tabular}
\end{center}
Baris $z$ ditulis sebagai $z-5x_{1}-4x_{2}=0$; koefisien negatif berarti menaikkan variabel
itu menaikkan $z$, jadi kolom dengan koefisien paling negatif masuk, di sini $x_{1}$ dengan
$-5$.

Uji nisbah menentukan siapa yang keluar: bagi ruas kanan dengan entri positif di kolom
masuk, lalu ambil nisbah terkecil. $24/6=4$ dan $6/1=6$, jadi $s_{1}$ keluar dan pivotnya
$6$. Membagi baris pivot dengan $6$ lalu menghapus kolom itu dari baris lain memberi
\[
  \begin{array}{rcl}
    x_{1}+\tfrac{2}{3}x_{2}+\tfrac{1}{6}s_{1} &=& 4,\\
    \tfrac{4}{3}x_{2}-\tfrac{1}{6}s_{1}+s_{2} &=& 2,\\
    z-\tfrac{2}{3}x_{2}+\tfrac{5}{6}s_{1} &=& 20 .
  \end{array}
\]
Titiknya kini $(4;0)$ berlaba $20$, dan masih ada koefisien negatif: $x_{2}$ dengan
$-2/3$. Uji nisbah memberi $4/(2/3)=6$ dan $2/(4/3)=1{,}5$, jadi $s_{2}$ keluar dan
pivotnya $4/3$. Sesudah satu putaran lagi,
\[
  \begin{array}{rcl}
    x_{1}+\tfrac{1}{4}s_{1}-\tfrac{1}{2}s_{2} &=& 3,\\
    x_{2}-\tfrac{1}{8}s_{1}+\tfrac{3}{4}s_{2} &=& \tfrac{3}{2},\\
    z+\tfrac{3}{4}s_{1}+\tfrac{1}{2}s_{2} &=& 21 .
  \end{array}
\]
Baris $z$ sekarang tidak punya koefisien negatif, jadi tidak ada arah yang menaikkan laba:
$(3;1{,}5)$ optimal dengan laba $21$. Seluruh aritmetika di atas diperiksa ulang dengan pecahan eksak, supaya tidak ada
pembulatan desimal yang menyelinap.

Baris terakhir itu lebih dari sekadar tanda berhenti. Koefisien $s_{1}$ dan
$s_{2}$ di sana, $3/4$ dan $1/2$, adalah harga dua bahan tadi. Angka itu muncul dari matriks
basis, $\mathbf{y}^{\top}=\mathbf{c}_{B}^{\top}B^{-1}$, dan
$\mathbf{b}^{\top}\mathbf{y}=24\cdot\tfrac34+6\cdot\tfrac12=21$, sama dengan laba.
Kesamaan itu adalah dualitas, dan bagian berikutnya menjelaskan mengapa ia sendiri sudah
merupakan bukti optimalitas.

Dua catatan. Aturan koefisien paling negatif tidak menjamin berhenti, sebab ada contoh kecil
dengan putaran tanpa akhir; aturan Bland yang memilih indeks terkecil menjamin berhenti.
Perangkat lunak modern memakai simpleks dual atau metode titik dalam, tetapi melaporkan
verteks dan harga yang sama.

\section{Dualitas: harga untuk setiap kendala}
\tingkat{lanjut}

Setiap bahan yang terbatas punya harga tersembunyi. Satu lembar kayu lapis tambahan
menaikkan laba, dan besar kenaikannya adalah harga lembar itu bagi bengkel. Harga itu tidak
tertulis pada soal aslinya, tetapi bisa dihitung, dan dualitas adalah cara menghitungnya.

Kalikan kendala kayu lapis dengan $y_{1}\ge0$ dan kendala jam dengan $y_{2}\ge0$, lalu
jumlahkan. Hasilnya
\[
  (6y_{1}+y_{2})x_{1}+(4y_{1}+2y_{2})x_{2}\le24y_{1}+6y_{2}.
\]
Bila $6y_{1}+y_{2}\ge5$ dan $4y_{1}+2y_{2}\ge4$, ruas kiri paling sedikit
$5x_{1}+4x_{2}$, yaitu laba. Jadi $24y_{1}+6y_{2}$ menjadi batas atas laba bagi setiap harga yang memenuhi dua syarat
itu, dan batas terbaiknya dibuat sekecil mungkin.

\begin{definisi}[Pasangan primal dan dual]
Untuk program linear
\[
  \text{(P)}\quad \text{maksimumkan } \mathbf{c}^{\top}\mathbf{x}
  \ \text{ dengan } A\mathbf{x}\le\mathbf{b},\ \mathbf{x}\ge\mathbf{0},
\]
\emph{program dualnya} adalah
\[
  \text{(D)}\quad \text{minimumkan } \mathbf{b}^{\top}\mathbf{y}
  \ \text{ dengan } A^{\top}\mathbf{y}\ge\mathbf{c},\ \mathbf{y}\ge\mathbf{0}.
\]
Variabel $\mathbf{y}$ disebut \emph{variabel dual} atau \emph{harga bayangan}.
\index{dualitas}\index{harga bayangan}
\end{definisi}

Untuk bengkel tadi, dualnya meminimumkan $24y_{1}+6y_{2}$ dengan kendala
$6y_{1}+y_{2}\ge5$, $4y_{1}+2y_{2}\ge4$, dan $y_{1},y_{2}\ge0$. Penyelesaiannya
$y_{1}=3/4$, $y_{2}=1/2$, bernilai $21$, sama dengan laba terbaik.

Bahwa batas atas dari dual tidak pernah di bawah nilai primal adalah hal pertama yang perlu
dipegang, dan buktinya tiga baris.

\begin{teorema}[Dualitas lemah]
Untuk setiap $\mathbf{x}$ layak (P) dan setiap $\mathbf{y}$ layak (D) berlaku
$\mathbf{c}^{\top}\mathbf{x}\le\mathbf{b}^{\top}\mathbf{y}$.
\index{dualitas lemah}
\end{teorema}

\begin{proof}
Karena $\mathbf{y}\ge\mathbf{0}$ dan $\mathbf{b}-A\mathbf{x}\ge\mathbf{0}$, berlaku
$\mathbf{b}^{\top}\mathbf{y}\ge\mathbf{y}^{\top}A\mathbf{x}$. Karena
$\mathbf{x}\ge\mathbf{0}$ dan $A^{\top}\mathbf{y}-\mathbf{c}\ge\mathbf{0}$, berlaku
$\mathbf{y}^{\top}A\mathbf{x}\ge\mathbf{c}^{\top}\mathbf{x}$. Rantai keduanya memberi
hasil yang diminta.
\end{proof}

Akibat praktisnya besar: sepasang titik layak, satu primal dan satu dual, dengan nilai sama
sudah membuktikan keduanya optimal tanpa memeriksa titik lain. Pada bengkel tadi,
$(3;1{,}5)$ dan $(0{,}75;\,0{,}5)$ memberi nilai $21$, jadi $21$ tidak bisa diperbaiki
maupun diturunkan.

Yang tidak diberikan dualitas lemah adalah jaminan bahwa kedua nilai itu selalu bertemu.
Jaminan itu isi teorema berikut.

\begin{teorema}[Dualitas kuat]
Bila (P) atau (D) mempunyai nilai optimum berhingga, maka keduanya mempunyai nilai optimum
dan nilainya sama.
\end{teorema}

Buktinya tidak dikerjakan di sini. Jalannya ada dua: lewat lemma Farkas tentang sistem
ketaksamaan linear, atau lewat kenyataan bahwa simpleks berhenti pada tabel tanpa koefisien
negatif, dan tabel itu sekaligus memberi $\mathbf{y}$ layak dual bernilai sama. Bukti
lengkapnya ada pada \texttt{WSPrinted.pdf} bagian dualitas program linear dan pada Boyd dan
Vandenberghe bagian 5.2.

Dari kesamaan itu lahir syarat yang dipakai untuk memeriksa hasil hitungan dengan cepat.
Kalimat sehari-harinya: kendala yang masih longgar tidak perlu dibayar, dan variabel yang
dipakai harus membayar harganya dengan pas.

\begin{akibat}[Kelengkapan komplementer]
Bila $\mathbf{x}$ dan $\mathbf{y}$ keduanya optimal, maka
\[
  y_{i}\bigl(b_{i}-\mathbf{a}_{i}^{\top}\mathbf{x}\bigr)=0 \ \text{ untuk setiap } i,
  \qquad
  x_{j}\bigl((A^{\top}\mathbf{y})_{j}-c_{j}\bigr)=0 \ \text{ untuk setiap } j .
\]
Sebaliknya, bila $\mathbf{x}$ layak (P), $\mathbf{y}$ layak (D), dan kedua kelompok
kesamaan itu berlaku, maka keduanya optimal.
\index{kelengkapan komplementer}
\end{akibat}

\begin{proof}
Karena keduanya optimal, $\mathbf{c}^{\top}\mathbf{x}=\mathbf{b}^{\top}\mathbf{y}$, dan
rantai pada bukti dualitas lemah menjadi rantai kesamaan. Setiap suku yang dijumlahkan di
dalamnya tak negatif, jadi setiap suku nol; suku-suku itu tepat
$y_{i}(b_{i}-\mathbf{a}_{i}^{\top}\mathbf{x})$ dan
$x_{j}((A^{\top}\mathbf{y})_{j}-c_{j})$. Arah sebaliknya: bila semua suku nol, rantai yang
sama memberi kesamaan nilai, dan menurut dualitas lemah nilai itu optimum kedua masalah.
\end{proof}

Sekarang harga bayangan itu sendiri. Bila ruas kanan diganggu, nilai optimum
berubah, dan laju perubahannya adalah komponen $\mathbf{y}$. Menambah satu lembar kayu lapis
menaikkan laba dari $21$ menjadi $21{,}75$; menambah satu jam perakitan menaikkannya menjadi
$21{,}5$. Jadi selembar kayu lapis bernilai $0{,}75$ satuan laba, satu jam perakitan
$0{,}5$.

Harga itu tidak berlaku selamanya, dan batas berlakunya bisa dihitung. Selama susunan
kendala aktif tidak berubah, $\mathbf{x}_{B}=B^{-1}\mathbf{b}$ linear dalam $\mathbf{b}$,
sehingga nilainya juga linear. Untuk bengkel tadi,
\[
  B^{-1}=\begin{pmatrix}\tfrac14&-\tfrac12\\[2pt]-\tfrac18&\tfrac34\end{pmatrix},
  \qquad
  \mathbf{x}_{B}=\begin{pmatrix}\tfrac14b_{1}-\tfrac12b_{2}\\[2pt]
  -\tfrac18b_{1}+\tfrac34b_{2}\end{pmatrix}.
\]
Dengan $b_{2}=6$ tetap, syarat $\mathbf{x}_{B}\ge\mathbf{0}$ memberi $b_{1}\in[12,36]$; di
luar selang itu satu variabel basis menjadi negatif dan kemiringannya berganti. Dengan $b_{1}=24$ tetap, selangnya $b_{2}\in[4,12]$; harga bayangan karena itu selalu
dilaporkan bersama selangnya, bukan sebagai angka telanjang.

Satu peringatan berguna sejak awal: harga bayangan adalah gagasan model bersambung. Pada
model integer, satu lembar kayu lapis tambahan bisa menaikkan laba sebesar $1$ atau tidak
sama sekali, dan tidak ada harga tunggal yang merangkum keduanya.

\section{Relaksasi LP: batas yang sah, pembulatan yang menipu}
\tingkat{menengah}

Kalau rak dan bangku boleh pecahan, soal tadi selesai: $3$ rak dan $1{,}5$ bangku.
Kenyataannya tidak, dan hubungan antara soal bersambung dan soal bulat perlu dinyatakan
dengan hati-hati.

\begin{definisi}[Program integer, relaksasi, dan celah integralitas]
\emph{Program integer} adalah program linear yang sebagian atau seluruh variabelnya
dituntut bulat. Bila seluruh variabel dituntut bulat, ia disebut program integer murni.
\emph{Relaksasi LP}-nya adalah program linear yang sama tanpa tuntutan kebulatan. Untuk
masalah maksimisasi dengan nilai optimum $z^{*}_{\mathrm{LP}}$ dan $z^{*}_{\mathrm{IP}}$,
\emph{celah integralitas} adalah $z^{*}_{\mathrm{LP}}-z^{*}_{\mathrm{IP}}$.
\index{program integer}\index{relaksasi}\index{celah integralitas}
\end{definisi}

\kamus{celah integralitas}{integrality gap}{Jarak antara batas yang dijanjikan model
bersambung dan jawaban bulat sebenarnya; jarak itu mengukur seberapa jauh model bersambung
menyesatkan.}

Alasan relaksasi berguna adalah satu kalimat: himpunan titik bulat yang layak adalah
bagian dari himpunan titik bersambung yang layak, jadi memaksimumkan di himpunan yang lebih
besar tidak mungkin memberi nilai yang lebih kecil.

\begin{proposisi}[Relaksasi memberi batas yang sah]
Untuk masalah maksimisasi, $z^{*}_{\mathrm{IP}}\le z^{*}_{\mathrm{LP}}$. Untuk masalah
minimisasi, ketaksamaannya berbalik.
\index{batas relaksasi}
\end{proposisi}

\begin{proof}
Setiap titik bulat yang layak bagi program integer juga layak bagi relaksasinya, sebab
relaksasi hanya menghapus syarat, tidak menambah syarat. Jadi daerah layak program integer
termasuk di dalam daerah layak relaksasi. Maksimum atas himpunan yang lebih besar tidak
lebih kecil daripada maksimum atas himpunan bagiannya.
\end{proof}

Batas itu sah, tetapi tidak selalu rapat, dan selisihnya bisa sangat menyesatkan bila
dipakai untuk membulatkan. Ada dua cara pembulatan gagal, dan keduanya muncul pada contoh
kecil kita.

Cara pertama: hasil pembulatan tidak layak. Titik $(3;1{,}5)$ dibulatkan ke atas menjadi
$(3;2)$. Labanya $23$, lebih tinggi daripada $21$, tetapi kayu lapisnya
$6\cdot3+4\cdot2=26$ lembar, melebihi $24$: itu bukan rencana produksi, hanya deretan
angka.

Cara kedua, lebih halus: hasil pembulatan layak tetapi bukan yang terbaik. Pembulatan ke
bawah memberi $(3;1)$ berlaba $19$, padahal $(4;0)$ berlaba $20$ dan seluruhnya layak.
Membulatkan koordinat satu demi satu ke arah terdekat karena itu memberi jawaban lebih buruk
daripada titik bulat lain yang jaraknya lebih jauh dari $(3;1{,}5)$.

Kedua kegagalan itu akibat bentuk daerah layaknya. Himpunan titik bulat layak
tidak konveks: $(3;1)$ dan $(4;0)$ anggotanya, tetapi titik tengahnya $(3{,}5;\,0{,}5)$
bukan, sebab tidak bulat. Pada himpunan tak konveks, teorema optimum di verteks tidak
berlaku: tidak ada verteks yang dijamin memuat jawaban, dan tidak ada tetangga yang bisa
diandalkan.

\begin{figure}[!ht]
\centering
\begin{tikzpicture}
\begin{axis}[width=9.3cm,height=6.3cm,
  xlabel={$x_{1}$ (rak)},ylabel={$x_{2}$ (bangku)},
  xmin=-0.35,xmax=5.1,ymin=-0.35,ymax=3.9,
  axis lines=left,xtick={0,1,2,3,4,5},ytick={0,1,2,3},
  tick label style={font=\scriptsize},label style={font=\scriptsize}]
\addplot[forget plot,fill=layerteal!8,draw=layerteal,thick]
  coordinates {(0,0) (4,0) (3,1.5) (0,3)} \closedcycle;
\addplot[forget plot,only marks,mark=*,mark size=1.5pt,draw=layerindigo]
  coordinates {(0,0) (1,0) (2,0) (3,0) (4,0)
               (0,1) (1,1) (2,1) (3,1)
               (0,2) (1,2) (2,2) (0,3)};
\addplot[forget plot,only marks,mark=square*,mark size=2.4pt,draw=layerjade]
  coordinates {(4,0)};
\addplot[forget plot,only marks,mark=*,mark size=1.5pt,draw=layerplum]
  coordinates {(3,1)};
\addplot[forget plot,only marks,mark=o,mark size=3pt,thick,accent]
  coordinates {(3,1.5)};
\draw[densely dashed,layergrey] (axis cs:3,1) -- (axis cs:3,2.35);
\node[font=\scriptsize,anchor=west,layerjade] at (axis cs:4.1,-0.05) {$(4;0)$, $z=20$};
\node[font=\scriptsize,anchor=east,layerplum] at (axis cs:2.9,0.85) {$(3;1)$, $z=19$};
\node[font=\scriptsize,anchor=west,accent] at (axis cs:3.15,1.6) {$(3;1{,}5)$, $z=21$};
\node[font=\scriptsize,anchor=north,layergrey] at (axis cs:3,-0.05) {};
\end{axis}
\end{tikzpicture}
\caption{Kisi titik bulat versus titik relaksasi. Ada $13$ titik bulat layak, ditandai titik
penuh. Relaksasi memberi $(3;1{,}5)$ bernilai $21$, yang berdiri di tengah antara $(3;1)$
bernilai $19$ dan $(3;2)$ yang tidak layak sebab $6\cdot3+4\cdot2=26>24$. Jawaban bulat
terbaik $(4;0)$ bernilai $20$ justru berjarak lebih jauh dari titik relaksasi.}
\label{fig:ch22-kisi-relaksasi}
\end{figure}

Gambar~\ref{fig:ch22-kisi-relaksasi} memuat seluruhnya: tiga belas titik bulat layak, dan
satu titik relaksasi yang berdiri di antara dua titik. Yang di antara itu bukan titik bulat,
dan yang terdekat belum tentu yang terbaik.

Pembulatan juga tidak punya arti tunggal di dimensi banyak. Dengan dua variabel,
``titik bulat terdekat'' masih bisa dibayangkan; dengan seratus variabel, yang terdekat belum
tentu layak, dan yang layak belum tentu terbaik. Praktiknya bukan membulatkan, melainkan
memakai batas relaksasi untuk memangkas pencarian.

\begin{tantangan}
Maksimumkan $6x_{1}+5x_{2}$ dengan $2x_{1}+3x_{2}\le12$, $4x_{1}+x_{2}\le13$, dan
$x_{1},x_{2}\ge0$ bulat. Relaksasinya memberi $(2{,}7;\,2{,}2)$ bernilai $27{,}2$.
Pembulatan ke bilangan bulat terdekat, $(3;2)$, bernilai $28$: lebih tinggi daripada nilai
relaksasi mana pun yang masih layak. Tunjukkan mengapa $28$ tidak boleh dipercaya, lalu
tentukan jawaban bulat yang benar beserta nilainya.

Petunjuk 1: periksa kelayakan, bukan hanya nilai. Angka yang lebih tinggi bisa datang dari
titik yang bukan rencana.

Petunjuk 2: jepit $x_{1}$. Bilangan bulat terbesar yang masih layak adalah $3$; untuk
$x_{1}=3$, hitung batas $x_{2}$ dari kedua kendala dan ambil yang lebih kecil.

Kalau mentok, kembali ke bagian Relaksasi LP: batas yang sah, pembulatan yang menipu.
\end{tantangan}

\section{Mengapa program integer lebih sulit}
\tingkat{menengah}

Sampai di sini pertanyaannya wajar: kalau relaksasi hanya menambah satu syarat, mengapa
program integer tidak diselesaikan dengan menyelesaikan relaksasinya lalu membulatkan?
Jawaban singkatnya, kesulitannya bukan pada mencari jawaban, melainkan pada memastikan
bahwa jawaban itu yang terbaik.

Tiga perbedaan menjelaskannya.

Pertama, tidak ada sertifikat optimalitas yang sederhana. Pada program linear,
sepasang $\mathbf{x}$ dan $\mathbf{y}$ bernilai sama adalah bukti selesai yang bisa diperiksa
dalam waktu sebanding ukuran data. Pada program integer, harga dari relaksasi hanya memberi
batas, bukan titik bulat yang mencapainya: dual dari relaksasi memberi batas bawah, bukan
jawaban.

Kedua, selisih antara batas dan jawaban bisa besar. Pada selimut verteks graf
lengkap $n$ simpul, relaksasinya membolehkan $x_{v}=1/2$ untuk semua $v$ dengan nilai $n/2$,
sedangkan jawaban bulatnya $n-1$; nisbahnya $2(n-1)/n$ mendekati $2$. Perhitungan lengkapnya
ada pada Laboratorium komputasi dan pada Latihan 10.

Ketiga, biaya komputasinya berbeda kelas, dan jumlah kandidatnya tumbuh eksponensial. Program
linear selesai dalam waktu polinomial terhadap ukuran masukan; hasil Khachiyan lewat metode
elipsoid dan Karmarkar lewat metode titik dalam adalah dua tonggaknya. Versi keputusan
program integer, ``apakah ada titik bulat layak bernilai paling sedikit $K$'', termasuk
NP-lengkap, dan daftar dua puluh satu masalah Karp memuat knapsack, yang sendirinya punya
$2^{n}$ himpunan bagian. Tidak ada algoritme polinomial yang diketahui untuk program integer,
dan bila ada, banyak hal lain akan runtuh bersamanya.

\begin{untukanak}
Pernahkah kamu memasang ubin di lantai? Kamu mengukur panjang ruangan, lalu membagi dengan
ukuran ubin. Hasilnya jarang bilangan bulat. Kalau kamu bulatkan ke bawah, ada celah di
ujung lantai. Kalau kamu bulatkan ke atas, kamu memotong ubin dan membuang sisanya. Dan
kadang jawaban terbaik bukan salah satu dari keduanya, melainkan pola pemasangan yang sama
sekali lain: mulai dari tengah, atau memakai ubin potongan di kedua sisi.

Orang yang hanya melihat satu angka hasil pembagian tidak akan menemukan pola
itu. Ia butuh dua batas: paling banyak berapa ubin terpakai, dan paling sedikit berapa.
Dengan dua angka itu ia berhenti memikirkan pola yang jelas terlalu boros, dan memusatkan
perhatian pada yang masih mungkin menang.

Di bab ini kita akan tahu mengapa batas dari model bersambung selalu sah, mengapa
pembulatan bisa memberi jawaban yang layak tetapi bukan yang terbaik, dan bagaimana
komputer memakai batas itu untuk memangkas jutaan kemungkinan tanpa memeriksanya satu per
satu.
\end{untukanak}

\section{Percabangan dan pembatasan}
\tingkat{lanjut}

Kalau jawaban relaksasi pecahan, kita tidak membulatkannya, melainkan memecah soalnya
menjadi dua dan bertanya lagi dengan relaksasi pada setiap potongan. Setiap jawaban bulat
memenuhi salah satu dari dua kemungkinan yang saling meniadakan, $x_{2}\le1$ atau
$x_{2}\ge2$, jadi menyelesaikan keduanya sama dengan menyelesaikan soal aslinya.

Alurnya empat langkah. Selesaikan relaksasi pada satu simpul. Bila hasilnya
bulat, catat sebagai calon jawaban dan jangan lanjutkan simpul itu. Bila pecahan, pilih satu
variabel yang pecah, $x_{j}=v$, lalu buat dua anak: $x_{j}\le\lfloor v\rfloor$ dan
$x_{j}\ge\lceil v\rceil$. Sebelum memeriksa sebuah simpul, bandingkan batas relaksasinya
dengan calon terbaik; bila tidak lebih baik, buang simpul itu beserta seluruh keturunannya.
Berhenti bila tidak ada simpul tersisa.

Aturan pembuangan itu sah, dan alasannya satu baris.

\begin{proposisi}[Pemangkasan tidak membuang jawaban]
Misalkan $z_{\mathrm{LP}}(N)$ adalah nilai relaksasi pada simpul $N$ dan $z_{\mathrm{best}}$
calon terbaik yang sudah tercatat. Bila $z_{\mathrm{LP}}(N)\le z_{\mathrm{best}}$, maka
tidak ada titik bulat layak di $N$ yang lebih baik daripada $z_{\mathrm{best}}$.
\end{proposisi}

\begin{proof}
Setiap titik bulat layak di $N$ juga layak bagi relaksasi $N$, jadi nilainya paling banyak
$z_{\mathrm{LP}}(N)$ menurut proposisi batas relaksasi. Bila
$z_{\mathrm{LP}}(N)\le z_{\mathrm{best}}$, semua titik itu bernilai paling banyak
$z_{\mathrm{best}}$, sehingga tidak ada yang perlu diperiksa.
\end{proof}

Pohon untuk bengkel tadi kecil, dan seluruhnya bisa ditulis sebagai tabel.
\begin{center}
\begin{tabular}{llcc}
\hline
simpul & tambahan kendala & $\mathbf{x}$ relaksasi & $z_{\mathrm{LP}}$ \\
\hline
akar & -- & $(3;\,1{,}5)$ & $21$ \\
kiri & $x_{2}\le1$ & $(10/3;\,1)$ & $20{,}\overline{6}$ \\
kanan & $x_{2}\ge2$ & $(2;\,2)$ & $18$ \\
kiri-kiri & $x_{2}\le1,\ x_{1}\le3$ & $(3;\,1)$ & $19$ \\
kiri-kanan & $x_{2}\le1,\ x_{1}\ge4$ & $(4;\,0)$ & $20$ \\
\hline
\end{tabular}
\end{center}
Akar memberi $21$ dengan $x_{2}$ pecah. Percabangan pada $x_{2}$ memberi dua anak: kiri
$20{,}\overline{6}$ dengan $x_{1}=10/3$ masih pecah, kanan $18$ dan sudah bulat. Simpul kanan
dicatat sebagai calon. Pada simpul kiri, percabangan pada $x_{1}$ memberi $(3;1)$ bernilai
$19$ dan $(4;0)$ bernilai $20$; yang kedua memperbarui calon. Tidak ada simpul tersisa, dan
jawabannya $20$.

\begin{figure}[!ht]
\centering
\begin{tikzpicture}[
  simpul/.style={draw=layergrey,fill=layergrey!6,rounded corners=1.5pt,
    inner sep=3pt,font=\scriptsize,align=center,text width=3.0cm},
  pangkas/.style={draw=accent,fill=accent!7,rounded corners=1.5pt,
    inner sep=3pt,font=\scriptsize,align=center,text width=3.0cm},
  panah/.style={-{Stealth},thick,draw=layergrey}]
\node[simpul] (akar) at (0,0) {akar\\ $z_{\mathrm{LP}}=21$\\ $(3;\,1{,}5)$};
\node[simpul] (kiri) at (-3.6,-1.75) {$x_{2}\le1$\\ $z_{\mathrm{LP}}=20{,}\overline{6}$\\
  $(10/3;\,1)$};
\node[pangkas] (kanan) at (3.6,-1.75) {$x_{2}\ge2$\\ $z_{\mathrm{LP}}=18$\\ $(2;\,2)$ bulat};
\node[pangkas] (kk) at (-5.1,-3.6) {$x_{1}\le3$\\ $z_{\mathrm{LP}}=19$\\ $(3;\,1)$ bulat};
\node[simpul] (kn) at (-2.1,-3.6) {$x_{1}\ge4$\\ $z_{\mathrm{LP}}=20$\\ $(4;\,0)$ bulat};
\draw[panah] (akar.south) -- (kiri.north);
\draw[panah] (akar.south) -- (kanan.north);
\draw[panah] (kiri.south) -- (kk.north);
\draw[panah] (kiri.south) -- (kn.north);
\node[font=\scriptsize,accent,anchor=north] at (3.6,-2.55) {dipangkas: $18<20$};
\node[font=\scriptsize,accent,anchor=north] at (-5.1,-4.45) {dipangkas: $19<20$};
\end{tikzpicture}
\caption{Pohon percabangan dan pembatasan untuk bengkel: lima simpul. Akar memberi $z=21$
dengan $x_{2}$ pecah. Cabang $x_{2}\ge2$ berakhir di $(2;\,2)$ bernilai $18$; cabang
$x_{2}\le1$ bercabang lagi menjadi $(3;\,1)$ bernilai $19$ dan $(4;\,0)$ bernilai $20$.
Kedua simpul berlabel bulat yang bernilai di bawah $20$ berbingkai merah: batasnya sudah
tidak lebih baik daripada calon terbaik, jadi keduanya dipangkas tanpa diperiksa lagi.}
\label{fig:ch22-pohon-percabangan}
\end{figure}

Gambar~\ref{fig:ch22-pohon-percabangan} menggambar pohon yang sama. Dua simpul di sana
berbingkai merah, yaitu $(2;\,2)$ bernilai $18$ dan $(3;\,1)$ bernilai $19$: keduanya bulat,
tetapi batasnya di bawah calon terbaik $20$, jadi keduanya dibuang.

Simpul kanan yang bernilai $18$ tidak perlu diperiksa lebih jauh seandainya
ia pecahan: batas $18$ sudah di bawah calon $20$, jadi seluruh cabangnya bisa dibuang. Pada
contoh sekecil ini penghematannya tidak terasa; pada soal dengan ratusan variabel, aturan
itulah yang membedakan selesai dalam semenit dan tidak selesai sama sekali, sebab tanpa
batas yang baik percabangan dan pembatasan merosot menjadi $2^{n}$ simpul.

Dua hal perlu jujur disebut. Waktu terburuknya tetap eksponensial, dan tidak ada
jaminan pohonnya kecil. Pemilihan variabel bercabang dan urutan simpul juga sangat
memengaruhi ukuran pohon; perangkat lunak memakai heuristik, misalnya variabel yang paling
dekat setengah. Keduanya menjelaskan mengapa program integer bisa jauh lebih lambat daripada
program linear berukuran sama.

\section{Contoh hitung}
\tingkat{menengah}

Sekarang seluruh angka dikerjakan sampai akhir, sehingga setiap klaim di atas bisa
diperiksa ulang dengan tangan atau dengan Python.

\begin{resep}
Empat langkah menyelesaikan program linear dua variabel secara grafik:
\begin{enumerate}
  \item Ubah tiap kendala menjadi garis dengan mengganti tanda pertidaksamaan menjadi tanda
        sama dengan, lalu gambar garisnya lewat dua titik potong sumbu.
  \item Tentukan sisi yang layak dengan menguji satu titik, biasanya titik nol; arsir sisi
        yang salah supaya daerah layak tetap bersih.
  \item Potongkan setiap pasangan garis, buang titik potong yang tidak memenuhi seluruh
        kendala, dan yang tinggal itulah verteks.
  \item Hitung nilai sasaran di setiap verteks, lalu ambil yang terbesar untuk maksimisasi
        atau terkecil untuk minimisasi.
\end{enumerate}
Peringatan urutan: langkah 3 tidak boleh dikerjakan sebelum langkah 2. Titik potong yang
dihitung lebih dulu sering terletak di luar daerah layak dan nilainya justru lebih tinggi
daripada jawaban sebenarnya; angka bagus yang muncul sebelum kelayakan diperiksa sulit
dibuang.
\end{resep}

\emph{Pertama, program linear bengkel.} Daerah layaknya segiempat dengan empat verteks,
dan keempatnya diperiksa:
\begin{center}
\begin{tabular}{lcc}
\hline
verteks & kendala aktif & laba $5x_{1}+4x_{2}$ \\
\hline
$(0;\,0)$ & $x_{1}=0,\ x_{2}=0$ & $0$ \\
$(4;\,0)$ & $6x_{1}+4x_{2}=24,\ x_{2}=0$ & $20$ \\
$(3;\,1{,}5)$ & $6x_{1}+4x_{2}=24,\ x_{1}+2x_{2}=6$ & $21$ \\
$(0;\,3)$ & $x_{1}+2x_{2}=6,\ x_{1}=0$ & $12$ \\
\hline
\end{tabular}
\end{center}
Titik $(3;1{,}5)$ dihitung dari kedua persamaan: $x_{1}=6-2x_{2}$ disubstitusi ke
$6x_{1}+4x_{2}=24$ memberi $36-8x_{2}=24$, jadi $x_{2}=3/2$, $x_{1}=3$, dan labanya
$15+6=21$.

\emph{Kedua, dual dan dualitas kuat.} Dualnya meminimumkan $24y_{1}+6y_{2}$ dengan
$6y_{1}+y_{2}\ge5$ dan $4y_{1}+2y_{2}\ge4$. Karena $x_{1},x_{2}>0$, kelengkapan
komplementer menuntut kedua kendala dual aktif; menyelesaikan keduanya sebagai kesamaan
memberi $y_{1}=3/4$ dan $y_{2}=1/2$. Nilainya $18+3=21$, sama dengan laba. Pemeriksaan silang:
$6\cdot\tfrac34+\tfrac12=5$ dan $4\cdot\tfrac34+2\cdot\tfrac12=4$, tepat sama dengan ruas
kanan, sehingga celahnya nol.

\emph{Ketiga, harga bayangan dan batas berlakunya.} Mengganggu ruas kanan memberi
tabel berikut, dengan $z$ dihitung ulang dari nol pada setiap baris.
\begin{center}
\begin{tabular}{lccc}
\hline
gangguan & $\mathbf{x}$ optimum & $z$ & perubahan \\
\hline
$b=(25;\,6)$ & $(3{,}25;\,1{,}375)$ & $21{,}75$ & $+0{,}75$ \\
$b=(24;\,7)$ & $(2{,}5;\,2{,}25)$ & $21{,}50$ & $+0{,}50$ \\
$b=(12;\,6)$ & $(0;\,3)$ & $12{,}00$ & $-9{,}00$ \\
$b=(24;\,4)$ & $(4;\,0)$ & $20{,}00$ & $-1{,}00$ \\
\hline
\end{tabular}
\end{center}
Baris ketiga mengurangi kayu lapis menjadi $12$, dengan perubahan
$-9=0{,}75\cdot(-12)$; baris keempat mengurangi $2$ jam, dengan perubahan
$-1=0{,}5\cdot(-2)$. Keduanya masih di dalam selang $b_{1}\in[12,36]$ dan
$b_{2}\in[4,12]$, jadi kelinierannya berlaku sampai ujung. Di luarnya kemiringan berganti:
dengan $b_{1}=40$, relaksasinya memberi $(0;\,3)$ bernilai $12$, bukan
$21+0{,}75\cdot16=33$.

\emph{Keempat, program integer bengkel.} Enumerasi dua arah tadi memberi jawaban $20$ di
$(4;0)$, sedangkan relaksasinya $21$ di $(3;1{,}5)$. Celahnya $1$, yaitu
$1/21=4{,}76\%$ dari nilai relaksasi. Pembulatan ke bawah memberi $19$, jadi kehilangan
$1$ satuan laba, atau $5\%$ dari jawaban sebenarnya, dan itu terjadi pada soal dengan dua
variabel saja.

\emph{Kelima, knapsack empat barang.} Empat barang dengan berat $5,4,3,2$ dan
nilai $10,7,5,3$ dimuat ke wadah berkapasitas $8$. Kerapatan nilainya $2$, $1{,}75$,
$1{,}667$, $1{,}5$, jadi aturan rakus mengambil barang pertama lalu $\tfrac34$ barang kedua:
berat $5+3=8$, nilai $15{,}25$. Jawaban bulat terbaiknya barang pertama dan ketiga, berat $8$
tepat penuh, nilai $15$. Celah integralitasnya $0{,}25$, hanya $1{,}64\%$ dari batas. Tetapi
pembulatan ke bawah membuang barang kedua sepenuhnya dan menyisakan nilai $10$, yaitu
$33{,}3\%$ di bawah jawaban sebenarnya: celah kecil tidak berarti pembulatan yang aman, sebab
keduanya mengukur hal berbeda.

\kertas{Untuk bengkel tadi, andaikan mesin perakitan hanya tersedia $5$ jam, bukan $6$.
Hitung ulang verteksnya, tentukan laba optimum bersambung, lalu tentukan laba optimum
bulatnya dengan enumerasi kecil. Sebelum menghitung, tulis dugaanmu: apakah celah
integralitasnya membesar atau mengecil? Sesudah itu periksa harga bayangan jam perakitan
pada $b_{2}=5$ dan bandingkan dengan $0{,}5$. Sebagai tambahan, cari $b_{2}$ tempat jawaban bulat terbaik berpindah dari $(4;0)$.}

\section{Lab komputasi}
\tingkat{lanjut}

Dua blok berikut dijalankan dengan Python 3.14.4, numpy 2.5.2, dan scipy 1.18.0. Blok
pertama menyelesaikan program linear, dualnya, dan gangguan ruas kanan. Blok kedua
menyelesaikan pohon percabangan, knapsack, dan selimut verteks pada graf lengkap lima
simpul. Keluaran disalin apa adanya.

\begin{lab}
\begin{verbatim}
import numpy as np
from scipy.optimize import linprog

np.set_printoptions(precision=6, suppress=True, linewidth=78)

c = np.array([5., 4.])              # laba per rak, per bangku
A = np.array([[6., 4.], [1., 2.]])  # kayu lapis, jam perakitan
b = np.array([24., 6.])

r = linprog(-c, A_ub=A, b_ub=b, bounds=[(0, None)]*2, method="highs")
print("LP   x* =", r.x, "  z* =", -r.fun)
print("sisa b - Ax =", b - A @ r.x)
print("marginals (scipy):", r.ineqlin.marginals)

rd = linprog(b, A_ub=-A.T, b_ub=-c, bounds=[(0, None)]*2, method="highs")
print("dual y* =", rd.x, "  nilai dual =", rd.fun)
print("celah dualitas =", rd.fun + r.fun)

print("B^-1 =", np.linalg.inv(A), "  x_B =", np.linalg.inv(A) @ b)

for nama, db in [("b=(25, 6)", (1., 0.)), ("b=(24, 7)", (0., 1.)),
                 ("b=(12, 6)", (-12., 0.)), ("b=(24, 4)", (0., -2.))]:
    rr = linprog(-c, A_ub=A, b_ub=b + np.array(db), bounds=[(0, None)]*2,
                 method="highs")
    print(f"{nama:<10} x = {np.round(rr.x, 6) + 0.0}  z = {-rr.fun:7.4f}"
          f"  delta z = {-rr.fun + r.fun:+.4f}")
\end{verbatim}
Keluaran sebenarnya:
\begin{verbatim}
LP   x* = [3.  1.5]   z* = 21.0
sisa b - Ax = [0. 0.]
marginals (scipy): [-0.75 -0.5 ]
dual y* = [0.75 0.5 ]   nilai dual = 21.0
celah dualitas = 0.0
B^-1 = [[ 0.25  -0.5  ]
 [-0.125  0.75 ]]   x_B = [3.  1.5]
b=(25, 6)  x = [3.25  1.375]  z = 21.7500  delta z = +0.7500
b=(24, 7)  x = [2.5  2.25]  z = 21.5000  delta z = +0.5000
b=(12, 6)  x = [0. 3.]  z = 12.0000  delta z = -9.0000
b=(24, 4)  x = [4. 0.]  z = 20.0000  delta z = -1.0000
\end{verbatim}
\end{lab}

\begin{lab}
\begin{verbatim}
import itertools
import numpy as np
from scipy.optimize import linprog, milp, LinearConstraint, Bounds

c = np.array([5., 4.])
A = np.array([[6., 4.], [1., 2.]])
b = np.array([24., 6.])


def lp_bengkel(tambahan_A=None, tambahan_b=None):
    Aub = A if tambahan_A is None else np.vstack([A, tambahan_A])
    bub = b if tambahan_b is None else np.concatenate([b, tambahan_b])
    r = linprog(-c, A_ub=Aub, b_ub=bub, bounds=[(0, None)]*2, method="highs")
    if r.status != 0:
        return (None, "tidak layak")
    return (-r.fun, np.round(r.x, 6) + 0.0)


print("-- pohon percabangan dan pembatasan --")
simpul = [("akar",        None, None),
          ("x2<=1",       [[0., 1.]],  [1.]),
          ("x2>=2",       [[0., -1.]], [-2.]),
          ("x2<=1,x1<=3", [[0., 1.], [1., 0.]],  [1., 3.]),
          ("x2<=1,x1>=4", [[0., 1.], [-1., 0.]], [1., -4.])]
for nama, ta, tb in simpul:
    z, x = lp_bengkel(None if ta is None else np.array(ta),
                      None if tb is None else np.array(tb))
    print(f"{nama:<14} z_LP = {z:8.4f}  x = {x}")

rb = milp(-c, constraints=LinearConstraint(A, -np.inf, b),
          integrality=[1, 1], bounds=Bounds(0, np.inf))
print("IP optimum   z_IP =", -rb.fun, " x =", np.round(rb.x, 6) + 0.0,
      " celah =", 21.0 + rb.fun)

print("-- knapsack 4 barang, kapasitas 8 --")
w = np.array([5., 4., 3., 2.])
v = np.array([10., 7., 5., 3.])
for mask in range(16):
    S = [i + 1 for i in range(4) if mask >> i & 1]
    bw, bv = w[[i - 1 for i in S]].sum(), v[[i - 1 for i in S]].sum()
    print(f"{str(S):<12} berat {bw:4.0f}  nilai {bv:4.0f}"
          f"  {'layak' if bw <= 8 else 'tidak layak'}")
rk = milp(-v, constraints=LinearConstraint(w, -np.inf, 8.),
          integrality=[1]*4, bounds=Bounds(0, 1))
print("enumerasi/MILP  z_IP =", -rk.fun, " x =", np.round(rk.x, 6) + 0.0)
rl = linprog(-v, A_ub=[w], b_ub=[8.], bounds=[(0, 1)]*4, method="highs")
print("relaksasi       z_LP =", -rl.fun, " x =", np.round(rl.x, 6) + 0.0)
print("celah integralitas =", -rl.fun + rk.fun)

print("-- selimut verteks K5: LP vs IP --")
n = 5
E = list(itertools.combinations(range(n), 2))
M = np.zeros((len(E), n))
for k, (i, j) in enumerate(E):
    M[k, i] = M[k, j] = 1.
rv = linprog(np.ones(n), A_ub=-M, b_ub=-np.ones(len(E)),
             bounds=[(0, 1)]*n, method="highs")
rc = milp(np.ones(n), constraints=LinearConstraint(M, 1, np.inf),
          integrality=[1]*n, bounds=Bounds(0, 1))
print("LP =", rv.fun, " x =", np.round(rv.x, 6) + 0.0)
print("IP =", rc.fun, " x =", np.round(rc.x, 6) + 0.0)
\end{verbatim}
Keluaran sebenarnya:
\begin{verbatim}
-- pohon percabangan dan pembatasan --
akar           z_LP =  21.0000  x = [3.  1.5]
x2<=1          z_LP =  20.6667  x = [3.333333 1.      ]
x2>=2          z_LP =  18.0000  x = [2. 2.]
x2<=1,x1<=3    z_LP =  19.0000  x = [3. 1.]
x2<=1,x1>=4    z_LP =  20.0000  x = [4. 0.]
IP optimum   z_IP = 20.0  x = [4. 0.]  celah = 1.0
-- knapsack 4 barang, kapasitas 8 --
[]           berat    0  nilai    0  layak
[1]          berat    5  nilai   10  layak
[2]          berat    4  nilai    7  layak
[1, 2]       berat    9  nilai   17  tidak layak
[3]          berat    3  nilai    5  layak
[1, 3]       berat    8  nilai   15  layak
[2, 3]       berat    7  nilai   12  layak
[1, 2, 3]    berat   12  nilai   22  tidak layak
[4]          berat    2  nilai    3  layak
[1, 4]       berat    7  nilai   13  layak
[2, 4]       berat    6  nilai   10  layak
[1, 2, 4]    berat   11  nilai   20  tidak layak
[3, 4]       berat    5  nilai    8  layak
[1, 3, 4]    berat   10  nilai   18  tidak layak
[2, 3, 4]    berat    9  nilai   15  tidak layak
[1, 2, 3, 4] berat   14  nilai   25  tidak layak
enumerasi/MILP  z_IP = 15.0  x = [1. 0. 1. 0.]
relaksasi       z_LP = 15.25  x = [1.   0.75 0.   0.  ]
celah integralitas = 0.25
-- selimut verteks K5: LP vs IP --
LP = 2.5  x = [0.5 0.5 0.5 0.5 0.5]
IP = 4.0  x = [1. 1. 1. 1. 0.]
\end{verbatim}
\end{lab}

Tiga hal perlu dibaca dari keluaran pertama. $\mathbf{x}^{*}$ yang dilaporkan
$(3;\,1{,}5)$ dengan $z^{*}$ tepat $21$, tanpa sisa pada kedua kendala, sesuai kelengkapan
komplementer. \texttt{scipy} melaporkan pengali $(-0{,}75,\,-0{,}5)$, negatif dari harga
bayangan, sebab rutin meminimumkan $-\mathbf{c}^{\top}\mathbf{x}$; angka benarnya $0{,}75$
dan $0{,}5$, dan cara memastikannya adalah menghitung ulang lewat $B^{-1}$. Celah dualitas
keluar $0$ tepat, bukan $10^{-15}$, karena $21$ terwakili tepat dalam biner; pada data
desimal, sisa orde $10^{-12}$ adalah hal biasa.

Keluaran kedua cocok dengan hitungan tangan: lima simpul dengan nilai yang sama
sampai empat desimal, jawaban bulat $20$ di $(4;\,0)$, relaksasi knapsack $15{,}25$, dan
jawaban bulatnya $15$ dengan celah $0{,}25$. Enumerasi keenam belas himpunan bagian
menegaskan $15$ yang terbaik: semua himpunan bernilai lebih tinggi melampaui kapasitas.

Bagian terakhir keluaran kedua memeriksa klaim faktor dua tadi. Untuk graf
lengkap lima simpul, relaksasinya $2{,}5$ sedangkan jawaban bulatnya $4$; nisbahnya $1{,}6$,
dan untuk $n$ simpul nisbahnya $2(n-1)/n$ menuju $2$ dari bawah. Batas relaksasi di sini bisa
meleset hampir dua kali.

\begin{untukdosen}
Urutan yang berhasil dipakai: mulai dari bengkel dan jawaban setengah bangku, lalu bentuk
baku, lalu verteks, baru simpleks. Jangan mulai dari tabel. Mahasiswa yang melihat tabel
lebih dulu menghafal langkah pivot tanpa tahu mengapa langkah itu boleh berhenti; yang sudah
melihat jawaban ada di sudut akan menganggap pivot sebagai cara berpindah sudut, dan uji
keoptimalan sebagai cara tahu tidak ada sudut yang lebih baik.

Dua miskonsepsi khas. Pertama, ``verteks adalah semua perpotongan garis
kendala''. Banyak perpotongan terletak di luar daerah layak, misalnya $(4;3)$ pada contoh
kita, yang memotong kedua garis tetapi melanggar keduanya; memeriksa kelayakan setiap kandidat
perlu dikerjakan di papan. Kedua, ``jawaban bulat adalah pembulatan jawaban pecahan''. Bengkel
memberi $19$ melawan $20$, knapsack memberi $10$ melawan $15$; keduanya cukup untuk
mematahkan dugaan itu.

Aktivitas kelas: tulis tabel awal di papan, minta kelas memilih kolom masuk dan menghitung
uji nisbah, lalu memprediksi laba sebelum pivot. Sesudah tabel akhir, tanyakan laba bila
kayu lapis menjadi $25$ lembar; $21{,}75$ membuka jalan ke selang $[12,36]$.

Tanda kesiapan lanjut: mahasiswa dapat menjelaskan mengapa metode simpleks tidak perlu
memeriksa seluruh verteks, dan dapat menyebut satu alasan mengapa relaksasi tidak boleh
dibulatkan.
\end{untukdosen}

\section{Ringkasan}

\begin{ringkasan}
\begin{itemize}
  \item Program linear menuntut sasaran dan kendala berbentuk garis. Bentuk baku
        $\max\mathbf{c}^{\top}\mathbf{x}$ dengan $A\mathbf{x}=\mathbf{b}$ dan
        $\mathbf{x}\ge\mathbf{0}$ berlaku untuk semua kasus, dengan variabel kendur untuk
        mengubah ketaksamaan menjadi kesamaan.
  \item Bila daerah layak terbatas dan sasarannya terbatas di atas, nilai terbaik dicapai di
        verteks. Pada bentuk baku, verteks sama dengan solusi layak basis, dan jumlahnya
        $\binom{n}{m}$, sehingga memeriksa semuanya bukan rencana masuk akal.
  \item Metode simpleks berjalan dari verteks ke verteks lewat kolom dengan biaya tereduksi
        negatif dan uji nisbah, dan berhenti ketika tidak ada biaya tereduksi negatif; saat
        itu baris $z$ sudah memuat harga dual.
  \item Dualitas lemah berlaku tanpa syarat; dualitas kuat menyatakan kedua nilai sama bila
        keduanya berhingga; kelengkapan komplementer menyatakan kendala yang longgar
        berharga nol.
  \item Harga bayangan $y_{i}$ adalah laju perubahan nilai optimum terhadap ruas kanan, dan
        hanya berlaku selama susunan kendala aktif tidak berubah; pada bengkel $0{,}75$ per
        lembar dan $0{,}5$ per jam, sah untuk $b_{1}\in[12,36]$ dan $b_{2}\in[4,12]$.
  \item Relaksasi LP memberi batas yang sah, bukan jawaban. Pembulatan bisa memberi titik
        tidak layak, seperti $(3;2)$, atau titik layak yang bukan terbaik, seperti $(3;1)$
        bernilai $19$ melawan $20$.
  \item Percabangan dan pembatasan memakai batas relaksasi untuk membuang cabang. Pohon
        bengkel berisi lima simpul, jawaban bulatnya $20$ dengan celah $1$. Waktu terburuknya
        tetap eksponensial, dan versi keputusannya NP-lengkap.
\end{itemize}
\end{ringkasan}

\section{Latihan}

\begin{latihan}
\begin{enumerate}
  \item Sebuah toko kue memakai $10$ kg tepung dan $8$ jam oven. Kue A memakai $2$ kg dan
        $1$ jam dengan laba $30$; kue B memakai $1$ kg dan $2$ jam dengan laba $20$.
        Tuliskan program linearnya, gambar daerah layaknya, dan hitung laba di setiap
        verteks. Lalu tentukan jawaban bulat terbaik dengan enumerasi, dan bandingkan
        dengan hasil pembulatan ke bawah dari jawaban bersambung.

  \item Buktikan bahwa himpunan $\{\mathbf{x}:A\mathbf{x}\le\mathbf{b}\}$ konveks, lalu
        jelaskan mengapa bukti itu tidak berlaku untuk himpunan
        $\{\mathbf{x}\in\mathbb{Z}^{n}:A\mathbf{x}\le\mathbf{b}\}$.

  \item Ganti laba bengkel menjadi $x_{1}+2x_{2}$. Hitung nilainya di keempat verteks,
        sebutkan verteks mana yang optimal, dan tunjukkan bahwa seluruh ruas di antara
        dua verteks optimal juga optimal. Terangkan mengapa titik di tengah ruas itu bukan
        verteks, lalu hubungkan hasilnya dengan harga bayangan kendala jam.

  \item Selesaikan dengan tabel simpleks: maksimumkan $3x_{1}+2x_{2}$ dengan
        $x_{1}+x_{2}\le4$, $x_{1}\le3$, $x_{2}\le3$, $x\ge0$. Tuliskan setiap tabel antara,
        sebutkan kolom masuk dan alasannya, lalu bacakan harga dualnya dari tabel akhir.
        (Petunjuk: dua pivot cukup; bandingkan hasil akhir dengan verteks yang diperoleh
        dari perpotongan $x_{1}+x_{2}=4$ dan $x_{1}=3$.)

  \item Tuliskan dual dari program pada latihan 1, selesaikan, dan periksa dualitas kuat
        secara numerik. Jelaskan arti setiap harga bayangan dalam satu kalimat.

  \item Nilai klaim ini: \emph{pembulatan ke bawah dari jawaban relaksasi selalu memberi
        titik bulat yang layak.} Tentukan benar, salah, atau belum cukup bukti, dan
        tunjukkan satu contoh penyangkal dengan angka.

  \item Nilai klaim ini: \emph{harga bayangan yang positif berarti kendala itu pasti
        menekan.} Tentukan benar, salah, atau belum cukup bukti, lalu jelaskan hubungan
        yang sebenarnya antara harga bayangan dan kelengkapan komplementer.
        (Petunjuk: periksa apa yang terjadi pada kendala dengan $y_{i}>0$ menurut
        $y_{i}(b_{i}-\mathbf{a}_{i}^{\top}\mathbf{x})=0$.)

  \item Untuk knapsack dengan berat $(5,4,3,2)$ dan nilai $(10,7,5,3)$ serta kapasitas
        $8$, hitung relaksasi LP-nya dengan tangan memakai urutan kerapatan, lalu bandingkan
        dengan hasil enumerasi. Hitung celah integralitasnya dan nyatakan sebagai persen
        dari batas. (Petunjuk: urutkan menurut nilai per satuan berat, lalu isi wadah
        sampai penuh; barang terakhir masuk sebagian.)

  \item Ulangi latihan 8 dengan kapasitas $7$. Tentukan relaksasi, jawaban bulat, dan
        celahnya, lalu jelaskan mengapa celahnya bisa berubah walaupun hanya satu angka
        data yang berubah.

  \item Pada selimut verteks graf lengkap $n$ simpul, buktikan bahwa relaksasi bernilai
        $n/2$ dan jawaban bulatnya $n-1$, lalu hitung nisbah keduanya untuk $n=4,5,10$.
        Jelaskan mengapa hasil itu menunjukkan bahwa pembulatan bukan cara umum untuk
        menyelesaikan program integer.
\end{enumerate}
\end{latihan}

\begin{untukpeneliti}
Yang benar-benar terbukti di bab ini: dualitas lemah untuk pasangan primal-dual dalam bentuk
baku, tanpa syarat apa pun; arah balik kelengkapan komplementer sebagai uji keoptimalan;
keoptimalan di verteks untuk daerah layak terbatas, dengan bukti lewat argumen kendala aktif;
dan kesahan pemangkasan pada percabangan dan pembatasan. Yang dipinjam: dualitas kuat, yang bukti
lazimnya lewat lemma Farkas atau lewat berhentinya simpleks; keduanya tercatat pada
\texttt{WSPrinted.pdf} bagian dualitas program linear dan pada Boyd dan Vandenberghe
bagian 5.2.

Yang dipopulerkan melampaui buktinya cukup banyak, dan tiga perlu disebut.
``Relaksasi LP sudah hampir benar'': faktor aproksimasi terburuknya bisa tumbuh tanpa batas
pada kelas masalah tertentu, dan bahkan pada selimut verteks mencapai $2$. ``Harga bayangan
adalah turunan nilai optimum'': itu benar hanya bila nilai optimum terdiferensialkan terhadap
ruas kanan; pada titik kemerosotan, saat dua basis optimal berbagi verteks yang sama,
turunannya tidak ada dan yang tersisa hanya subgradien. ``Cabang dan batas selalu cepat'':
waktu terburuknya eksponensial, dan kelambatan pada instans tertentu bukan tanda implementasi
yang buruk.

Pertanyaan terbuka: seberapa kuat hierarki pemotongan Chvatal-Gomory untuk kelas polytope
tertentu, dan apakah ada batas bawah superpolinomial untuk ukuran formulasi diperluas
polytope traveling salesman; hasil Fiorini dan kawan-kawan pada 2012 menunjukkan batas
eksponensial untuk kasus itu, sementara pencarian batas serupa untuk masalah lain
berjalan. Batas keberlakuan bab ini: data rasional, aritmetika eksak
untuk pembuktian, toleransi titik mengambang untuk praktik; seluruh pernyataan tentang harga
bayangan berlaku pada model bersambung, bukan model integer.

\peta{Bab ini mengikuti \emph{Introduction to Linear Optimization} pada berkas
\texttt{WSPrinted.pdf}: bagian formulasi, geometri himpunan polihedral (verteks dan titik
ekstrem), sistem ketaksamaan linear dan dualitas (kelengkapan komplementer, analisis
sensitivitas), serta simpleks (solusi layak basis, biaya tereduksi, uji keoptimalan). Sisi
program integer mengikuti Wolsey, \emph{Integer Programming}, bab 1 (formulasi) dan bab 2
(optimalitas, relaksasi, batas), yang menempatkan relaksasi sebagai sumber batas dan
percabangan sebagai cara memakainya. Sudut pandang konveks pada bagian kedua dan keempat
adalah Boyd dan Vandenberghe, \emph{Convex Optimization}, bab 4 dan 5. Hubungan dengan bab
lain: bentuk baku dan daerah layak melanjutkan Bab 21, penyelesaian sistem linear memakai
Bab 8, dan celah integralitas dipakai lagi pada Bab 25.}

\end{untukpeneliti}
